\documentclass[../main.tex]{subfiles}

\begin{document}
    \section{Introduction}

    \begin{definition}{Holomorfic function}{def:holomorfic}
        A function $f \colon \C \supseteq \Omega \to \C$ is \textit{holomorphic} 
        if it is complex differentiable on $\omega$, meaning that for any
        $z \in \Omega$
        \begin{equation*}
            f'(z) = \lim_{h \to 0} \frac{f(z + h) - f(z)}{h}
        \end{equation*}
        exists. Equivalently,
        \begin{equation*}
            f(z + h) = f(z) + h f'(z) + o(h).
        \end{equation*}
    \end{definition}

    Equivalently, $f$ is differentiable as a function $\R \to \R$ and its
    differential is a matrix of multiplication by a complex number
    \begin{equation*}
        \begin{bmatrix}
            a  & b \\
            -b & a
        \end{bmatrix}
    \end{equation*}
    for some reals $a$, $b$. Writing this using partials, we get

    \begin{definition}{Cauchy - Riemann equations}{def:cauchy-riemann}
        Let $f$ be holomorphic and write $f(x + iy) = u(x, y) + i v(x, y)$. Then
        \begin{equation*}
            \frac{\partial u}{\partial x} = \frac{\partial v}{\partial y}
            \text{ and }
            \frac{\partial v}{\partial x} = -\frac{\partial u}{\partial y}.
        \end{equation*}
        TODO: $\frac{\partial}{\partial z}$, $\frac{\partial}{\partial \overline{z}}$.
    \end{definition}
    
    \begin{lemma}{Closedness of $f dz$}{lem:closedness-dz}
        Let $dz = dx + i dy$. Then $f$ is holomorphic iff the form $f dz$ is closed, meaning
        \begin{equation*}
            d(f dz) = 0.
        \end{equation*}
    \end{lemma}

    All arithmetic properties are the same as for reals.

    \begin{example}{Some holomorphic functions}{exm:holomorfic:fns}
        Any polynomial, convergent power series (not proven yet), $\exp$,
        all trigonometric functions outside their poles. $\sin$ and $\cos$
        are unbounded for large $y$ and $\tan$ is bounded for large $y$.  
    \end{example}
    
\end{document}